
\Paragraph{Basics}
LSM-trees store data as key-value pairs. They have N levels and the first level, Level0, is stored in main memory and referred to as a buffer.
The remaining N-1 levels are stored on external storage.
RocksDB uses leveling instead of tiering which means that all levels consist of one run.(EXPLAIN MORE.)
These N-1 levels store data in the SST files.
SST stands for Sorted Sequence Table and they are persistent files storing data. 
In the file, keys are usually organized in sorted order so that a key or iterating position can be identified through a binary search.
All new entries are initially stored in the buffer and are later compacted to lower(HIGHER?) levels.
Compaction pushes keys and values to lower levels to free up space in memory.
There are many compaction policies, some more aggressive and others not so much.
This approach enables the system to treat updates, deletes and inserts very quickly as only the main memory is used.
If an update or delete on a key is requested, the system first searches in the main buffer by first consulting a dedicated bloom filter, and if the key exists, it is in-place updated or permanently deleted.
Otherwise, if the key is not in the buffer, instead of scanning the disk the system will insert a new entry in the tree.
This entry has absolutely no difference with a normal insert entry and will be eventually stored in an SST file. 
In the case of deleted the new entry has the same key as the one that needs to be deleted and a flag which showcases that it is a tombstone.
This is an out-of-place delete as the value is not immediately deleted.
Updates, work the same way, by inserting the newer version in memory but before that, a delete has to be initiated as the older version needs to be invalidated and later be deleted during compactions. 
The invalidated entries are eventually deleted when the file that contains the tombstone is compacted with the file that contains the value. 


\Paragraph{Partial Compaction} 
To amortize the latency spikes resulting from compactions in larger levels, state-of-the-art LSM based systems organize runs into smaller files, and perform compactions at the granularity of files instead of levels [22]. 
If Level i grows beyond a threshold capacity, a compaction is triggered and one file from Level i is chosen to be compacted with files from Level i + 1 that have an overlapping key-range, performing a partial compaction. 
Deciding which file to compact depends on the storage engine design. 
For instance, to optimize for high write throughput, we select the file from Level i that has minimal overlapping files in Level i + 1, to minimize both write amplification and compaction time. 
Partial compactions spread the total number of I/Os performed during compactions of an entire level over time to multiple smaller compactions.

\Paragraph{Point deletes in RocksDB}
Write something generic?? 

\Paragraph{Point deletes read path implementation in RocksDB} 
Deleting keys is a very common pattern in RocksDB.
A point delete  does not delete the actual value, but inserts a new entry to the LSM tree that is called "tombstone".
Every tombstone has its own sequence number.
The sequence number is an auto-incremented ID number that is used to implement snapshot read, garbage collection in compactions, MVCC in transactions, and some other purposes such as determining the age of the entries. 
The larger the sequence number, the newer the entry.
When a point lookup is requested from the database, the newest entry is retrieved. If this entry is a tombstone it means that the key is deleted. 

\Paragraph{Range deletes read path implementation in RocksDB} 
Most systems built on top of RocksDB have multi-component key schemas, where keys sharing a common prefix are logically related. 
The most common pattern is scan-and-delete, i.e., advance an iterator through the to-be-deleted range, and issue a delete for each key.
This approach, even though functional majorly increases space amplification as it creates many tombstones, slows down iterators and doesn't offer a deadline for space reclamation and increases I/Os. 
Not until recently, a dedicated mechanism for range deletes was implemented. 
However, even though it offers greater speeds in inserting delete ranges, it decreases look up speeds for non deleted keys. 
Also, the current system does not adapt and as a result executes `wasted' IOs in cases where a deleted key by a range delete is requested multiple times. 
Our method uses tries~\cite{Zhang2018a} in order to efficiently store keys that led to wasted IOs and avoid making them in the future. 
The more the system operates, the more useful our tries become, thus saving IOs which is the bound for most systems.

When a range delete is issued, it is first written to a dedicated memtable for range tombstones.
As memtable we refer to the in-memory data structure that stores the most recent updates of the database.

\Paragraph{Range tombstones} 
A range tombstone is composed of a start key, end key, and sequence number. 
Any key that falls within the range is considered deleted if the key's sequence number is less than or equal to the range tombstone's sequence number. 
Range tombstones are not fragmented in the memtable directly, but are instead fragmented each time a read occurs.
Range tombstones are stored in the SST files themselves. Not in the data blocks, but in a separate meta block for range tombstones. 
This resolves the problem of when to obsolete range tombstones, as it's simple: when they're compacted to the bottom level. 
Each SST file contains all range tombstones at that level that cover user keys overlapping with the file's key range; this greatly simplifies iterator seeking and point lookups. 
When a new range delete is requested from the system, a range delete tombstone is created which is later applied to a dedicated range tombstone memtable. 
Later in the background the range tombstone memtable and its corresponding data memtable are flushed together into a single SST with a range tombstone meta-block. 
As new data is inserted, new compactions are executed. 
During these compactions, in the same way as with single key tombstones, obsolete range delete tombstones are permanently deleted, thus reducing storage occupancy. 

\Paragraph{Range deletes read path Implementation in RocksDB} 
In a point lookup, RocksDB first looks at the fence pointers table kept in memory. 
This table keeps the page address and the first key of each page of the SST files. 
Then, by using binary search, the corresponding page is found. 
Afterwards, the bloom filter of each file probabilistically responds if the key is contained in the file. If not the key does not exist for sure. 
If yes, we execute an I/O to fetch the key and its sequence number. The next step is to determine if this key is deleted. 
If the key is deleted by a point delete it means that the most recent version of that key in the LSM-tree is a tombstone, which is practically an entry with a flag to indicate that it is a delete. 
In a range delete however, when the requested key resides in a range deleted area, the fence pointer and bloom filter will bypass this range tombstone and lead to the SST file that the key is stored. 
Thus, there is no immediate way to decide whether the key is deleted.

\Paragraph{Skyline vs RDA} 
In the the initial implementation of range deletes in RocksDB, each read operation  of built up a 'skyline' for the range tombstones visible at the read sequence number. 
In point lookups, range tombstones were aggregated in an unordered vector, essentially the 'skyline' as they were search through live memtable, immutable memtables, and then SSTs. 
This Skyline approach can very quickly determine whether a key is deleted but is expensive to construct and maintain. 
The initial skyline range tombstone implementation has since been replaced with a more efficient lookup structure. 
The key change in the new implementation is to fragment the range tombstones that are stored in an sstable. 
The fragmented range tombstones provide the same benefit as the skyline representation: the ability to binary search the fragments in order to find the tombstone covering a key. 
But unlike the skyline approach, the fragmented tombstones can be cached on a per-sstable basis. 
In the v2 approach, a range delete aggregator (RDA) keeps track of the fragmented range tombstones for each sstable encountered during a read or iterator, and logically merges them together. 
In summary, in RocksDB the range delete aggregator acts as an oracle for answering whether a key is deleted at a particular sequence number. 
Due to caching of fragmented tombstones, the v2 implementation of RDA is significantly faster to populate than v1, yet the overall approach to processing range tombstones remains similar. 
The entire path is visible in figure 2

\begin{figure}[h!]
    \centering
    \includegraphics[scale=0.3]{figures/read_path.png}
    \caption{Read path.}
    \label{fig:method}
\end{figure}
